\section{Discussion}
In this paper we discussed how disentangled representation learning aligns with the goals of subgroup fair machine learning, and presented a method for learning a structured latent code using multiple sensitive attributes.
The proposed model, FFVAE, provides \textit{flexibly fair} representations, which can be modified simply and compositionally at test time to yield a fair representation with respect to multiple sensitive attributes and their conjunctions, even when test-time sensitive attribute labels are unavailable.
Empirically we found that FFVAE disentangled sensitive sources of variation in synthetic image data, even in the challenging scenario where attributes and labels correlate.
Our method compared favorably with baseline disentanglement algorithms on downstream fair classifications by achieving better parity for a given accuracy budget across several group and subgroup definitions.
FFVAE also performed well on the Communities \& Crime and Celeb-A dataset, although none of the models performed robustly across all possible subgroups in the real-data setting.
This result reflects the difficulty of subgroup fair representation learning and motivates further work on this topic.

There are two main directions of interest for future work.
First is the question of fairness metrics: a wide range of fairness metrics beyond demographic parity have been proposed \citep{hardt2016equality,pleiss2017fairness}.
Understanding how to learn flexibly fair representations with respect to other metrics is an important step in extending our approach.
Secondly, robustness to distributional shift presents an important challenge in the context of both disentanglement and fairness.
In disentanglement, we aim to learn independent factors of variation.
Most empirical work on evaluating disentanglement has used synthetic data with uniformly distributed factors of variation, but this setting is unrealistic. 
Meanwhile, in fairness, we hope to learn from potentially biased data distributions, which may suffer from both undersampling and systemic historical discrimination.
We might wish to imagine hypothetical ``unbiased'' data or compute robustly fair representations, but must do so given the data at hand.
While learning fair or disentangled representations from real data remains a challenge in practice, we hope that this investigation serves as a first step towards understanding and leveraging the relationship between the two areas.
